\section{Response-Aware Marginal-Utility Routing}

R-MUR estimates marginal utility with the fixed predictor. Every candidate is
first represented by cached context features. At state $A$, let
$\widehat m_C(j\mid A)$ denote a cached-only utility score. A shortlist is
\begin{equation}
 \mathcal Q_A=\operatorname{TopQ}_{j\notin A}\widehat m_C(j\mid A),
 \qquad |\mathcal Q_A|=q\ll K.
\end{equation}
The shortlist reduces the number of predictor evaluations while retaining the
candidates that are most promising under the cached representation.

For each retained candidate, R-MUR computes a compact response summary from
the fixed predictor,
\begin{equation}
 r_{A,j}
 =
 \psi\!\left(f_A(x),f_{A\cup\{j\}}(x)\right).
 \label{eq:response}
\end{equation}
The summary contains the level of the current prediction, the level after
adding the candidate, the mean and standard deviation of the response change,
its mean absolute value, its normalized Euclidean norm, and its maximum
absolute value. These statistics describe how the candidate changes the
predictor, not merely how similar the candidate is to the query.

The response-aware utility model scores retained candidates with
\begin{equation}
 \widehat m_R(j\mid A)
 =
 g(h_Q,h_A,h_j,r_{A,j}),
 \qquad j\in\mathcal Q_A,
 \label{eq:rmur}
\end{equation}
where $h_Q$, $h_A$, and $h_j$ are cached representations of the query, the
selected set, and the candidate. The set encoder is permutation invariant.
Training minimizes
\begin{equation}
 \mathcal L
 =
 \mathcal L_{\mathrm{Huber}}
 +
 \beta
 \sum_{i,j:\,\widetilde m_i>\widetilde m_j}
 |\widetilde m_i-\widetilde m_j|
 \log\!\left(
 1+\exp\!\left[-(\widehat m_i-\widehat m_j)\right]
 \right),
 \label{eq:loss}
\end{equation}
where the pairwise term is evaluated within the same selection state. The gap
weight gives larger influence to pairs that can change the argmax decision.
The regression term keeps the utility scale calibrated on states without a
dominant alternative.

At inference, R-MUR starts from $A_0=\emptyset$ and repeats the following
steps until the budget is reached or a stopping condition is met. It scores
all remaining candidates with the cached model, forms $\mathcal Q_A$, computes
response summaries for the shortlist, reranks the shortlisted candidates with
$\widehat m_R$, and adds the highest-scoring positive candidate. The selected
set is updated and the process repeats. The procedure is greedy; it targets
one-step decision quality under a context budget.

A calibration-derived threshold controls when routing stops. On a held-out
split, let $r_i$ be the absolute residual between predicted and observed
marginals. For level $\alpha$, the radius
$q_\alpha=\operatorname{Quantile}_{1-\alpha}(\{r_i\})$ defines a symmetric
interval
\begin{equation}
 [L_{A,j},U_{A,j}]
 =
 [\widehat m_R(j\mid A)-q_\alpha,\,
  \widehat m_R(j\mid A)+q_\alpha].
\end{equation}
This threshold changes only the stopping decision; the candidate ranking
remains the response-aware score.

\subsection{A head shaped by the exact identity}
\label{sec:bilinear-head}

The experiments in \cref{sec:strong-backbone} show that a generic scalar head
over concatenated state, candidate and response features does not recover
marginal utility when the fixed predictor is strong. The analysis in
\cref{sec:analysis} says why, and prescribes the functional form. Conditional
on the state, \cref{eq:prop4} gives

\begin{equation}
 \E\!\left[\widetilde{\marginal}(j\mid A)\mid x,A\right]
 =
 \langle g_A(x), \dAj(x)\rangle-\|\dAj(x)\|^2-\lambda c_j ,
 \qquad g_A(x)=2(\mu(x)-f_A(x)).
\end{equation}

Only $g_A$ has to be learned: the response $\dAj$ is observed once the
predictor has been probed, and the curvature penalty is an exact function of
that same response. We therefore replace the scalar head by a \emph{bilinear
response head} that emits a vector in prediction space and scores

\begin{equation}
 \widehat m_B(j\mid A)
 =
 \langle g_\theta(\phi_A(x), c_j),\, \dAj(x)\rangle
 -\|\dAj(x)\|^2 - \lambda c_j ,
 \label{eq:bilinear-head}
\end{equation}
where $\phi_A(x)$ collects the same state and candidate features as before and
$g_\theta$ is a small network. This is the minimal-capacity model consistent
with the identity: it cannot represent a score that ignores the response, it
never has to learn the curvature term, and the only statistical task left is
estimating the state-dependent coefficient. The training target is unchanged
(the realised marginal), so the head is trained with the same regression and
gap-weighted ranking objective as before; the difference is purely inductive
bias. For a general smooth loss the same construction uses the identity of
\cref{prop:bregman} in place of the squared-loss form, and the cross-entropy
bound of \cref{prop:ce} replaces the exact penalty by its upper bound.

